\section{Re-evaluation}

Chaque DAG d'évaluation représente l'histoire de l'évolution d'une orbite dans
le contexte d'une évaluation initiale. %
Dans le cadre d'une réévaluation, l'édition de la spécification paramétrique
fait que l'histoire d'un paramètre topologique est susceptible d'être changée. %
Il est alors nécessaire d'adapter les DAGs d'évaluation à ce nouveau processus
de modélisation et ainsi permettre l'appariement des paramètre topologique. %
À l'issue de l'appariement, un DAG de réévaluation désignera alors un ou
plusieurs paramètres topologiques sur lesquels appliquer les opérations de
modélisation. %

\subsection{Re-evaluation DAG}

\begin{figure}[ht]
    \centering
    \includegraphics[height=.3\textheight]{pn3-reval}
    \caption{$PN_3$ re-evaluation DAG}
    \label{fig:pn3-reval-dag}
\end{figure}


Dans un DAG de réévaluation, les noms de nœuds sont remplacés par des noms
concrets de brins. %
Cela s'explique par le fait qu'un DAG de réévaluation est construit tout au long
du processus de réévaluation d'un objet représentant alors, concrètement,
l'histoire d'une orbite au fur et à mesure que l'objet est transformé par les
opérations appliquées successivement. %

Considérons l'exemple du DAG représentant l'histoire réévaluée du $PN_{3}$
(Fig.~\ref{fig:pn3-reval-dag}). %
As a remainder (Pour rappel), une application \texttt{ADD1-insert} est ajoutée
entre les application \texttt{1-square} et \texttt{2-extrusion} et l'application
\texttt{3-insert} est supprimée dans la spécification paramétrique. %
\begin{description}
  \item[Level 1] The application \texttt{1-square} creates the volume face $\orbit{0,1}(5)$.
  \item[Level 2] The added application \texttt{ADD1-insert} does not modify the volume face $\orbit{0,1}(5)$.
  \item[Level 3] From this volume face, \texttt{2-extrusion} creates the volume face $\orbit{0,1}(47)$.
  \item[Level 4] Finally, this latter volume face is actually not modified by \texttt{3-insert} application which is deleted in the re-evaluation.
\end{description}

Puisque l'ajout, la suppression ou encore le déplacement d'opération dans une
spécification paramétrique éditée peuvent provoquer des changements topologiques
dans le modèle réévalué. %
Il apparaît logique que ces éditions impactent les paramètres topologiques de
telle sorte qu'ils puissent être dupliqués voire supprimés. %

\begin{figure}[ht]
    \centering
    \includegraphics[width=\columnwidth]{pn7-reval}
    \caption{$PN_7$ re-evaluation DAG}
    \label{fig:pn7-reval-dag}
\end{figure}

Dans le premier cas, un paramètre topologique peut être dupliqué par l'ajout
d'une application provoquant la scission d'une entité topologique. %
Une stratégie pourrait être, pour une opération référençant un tel paramètre
topologique, sur chacun des paramètres dupliqués par la scission.
Par exemple, dans le jeu initial (Fig.~\ref{fig:evaluation}), l'application
\texttt{5-triangulation} ne triangule qu'une seule face de volume
($\orbit{0,1}$-orbit) (Fig.~\ref{fig:5-eval}). %
Or, l'ajout de l'application \texttt{ADD1-insert} dans la réévaluation
(Fig.~\ref{fig:reevaluation}), scinde l'arête de face de volume
($\orbit{0}$-orbit) et deux faces sont alors créées par l'application de
\texttt{2-extrusion} au lieu d'une. %
Le DAG de réévaluation du $PN_{7}$ illustre ce cas de figure, l'application
\texttt{ADD1-insert} au niveau 2 ajoute une scission décrivant deux orbites
faces issues de la même arête scindée mais incidentes à deux brins différents. %
Ce DAG propose, comme stratégie, d'appliquer \texttt{5-triangulation} sur les
faces incidentes aux brins 37 et 52. %

Pour la suppression d'un paramètre topologique, il existe plusieurs cas de
figures. %
Si l'orbite incidente brin désigné comme paramètre topologique est supprimée, il
nécessaire d'annuler l'application de toute opération référençant ce paramètre
topologique car celles-ci échoueront. %
Autrement, si l'orbite n'est pas entièrement supprimée, il est possible de
désigner arbitrairement un autre brin de ses brins afin de poursuivre la
réévaluation. %

\subsection{Parameter Matching}

\begin{figure}[ht]
    \centering
    \includegraphics[width=\columnwidth]{pn6-reval}
    \caption{$PN_6$ re-evaluation DAG}
    \label{fig:pn6-reval-dag}
\end{figure}

Un DAG de réévaluation se distingue d'un DAG d'évaluation par, entre autres,
l'usage de noms concrets au lieu de noms de nœuds. %
Les noms concrets sont obtenus par appariement des paramètres topologiques lors
de l'étape de réévaluation d'une spécification paramétrique. %

Reprenons les DAGs d'évaluation et de réévaluation du $PN_{3}$. %
Au premier niveau du DAG d'évaluation, le niveau d'orbite désigne une orbite
incidente au nœud $n_{5}$. %
Lors de la construction du DAG de réévaluation, il faut désigner un brin parmi
ceux filtrés par ce nœud. %
Puisque $5$ est le seul brin créé par le nœud $n_{5}$ de \texttt{1-square},
celui-ci est utilisé. %

%%% Local Variables:
%%% mode: latex
%%% TeX-master: "../main"
%%% End:
